\chapter{Literature Review}\label{chap:litreview}

%% ============================================================
%% PART I — Automated ICD Coding (from Old Report)
%% ============================================================
\section{Part I: Automated ICD Coding}\label{sec:p1:icd-lit}

\subsection{Introduction and the Evolution of ICD Coding}\label{sec:p1:icd-evolution}

The automation of the International Classification of Diseases (ICD) coding process has long been recognized as a critical area of research in medical informatics and Natural Language Processing (NLP). The ICD system, maintained by the World Health Organization, is used worldwide to classify diseases and health conditions for epidemiological research, healthcare management, and clinical billing. Accurate and efficient coding is essential because ICD codes capture details of a patient's diagnosis, comorbidities, and complications, thereby supporting reliable record-keeping and reimbursement processes~\citep{who2019icd11}.

Originally, the Ninth Revision (ICD-9) provided a limited number of disease categories. Over time, ICD-9 was phased out in favor of ICD-10, driven by the need for more specificity and additional capacity to accommodate new diseases~\citep{dong2022automated}. ICD-10-CM (Clinical Modification) expanded the code set substantially (from approximately 13,000 to 68,000 codes) and introduced alphanumeric codes capable of encoding disease severity, anatomical site, and laterality~\citep{edin2023automated}. Although this granularity supports more detailed patient records, it also increases complexity in manual coding workflows~\citep{farkas2008automatic}.

The progression to ICD-11~\citep{who2019icd11} further reflects international efforts to enhance coding flexibility and align medical concepts more precisely with emerging public health priorities. Its hierarchical structure and additional extensions aim to capture nuances of disease presentation, enabling more robust epidemiological tracking. However, these regular revisions amplify the complexity for human coders, who must keep abreast of ever-expanding vocabularies and specialized rules~\citep{scheurwegs2017data}. The high stakes of accurate billing and statistical reporting underscore the importance of developing automated coding solutions that can function reliably, even when new codes and terminologies are introduced.

This section reviews the broad spectrum of approaches and challenges that characterize automated ICD coding. Early rule-based and machine learning (ML) methods are summarized, followed by deep learning approaches that leverage convolutional, recurrent, and transformer-based architectures. Recent improvements in long-context transformer designs, modern BERT-like encoders, knowledge-based methods, and curriculum learning strategies are discussed, highlighting progress in handling the extensive and imbalanced ICD label space. The section also explores data imbalance and rare-code prediction, synthetic data generation for underrepresented diagnoses, interpretability, and the practical challenges of deploying automated ICD coding systems in real-world clinical environments. Relevant evaluation metrics and future research directions are presented to offer a holistic overview of current capabilities and ongoing developments in this domain.

\subsection{Early Approaches to Automated ICD Coding}\label{sec:p1:early-approaches}

\subsubsection{Rule-Based Systems and Traditional Machine Learning}

Early efforts to automate ICD coding primarily involved rule-based systems and traditional machine learning techniques. Rule-based systems utilized handcrafted linguistic rules, domain-specific ontologies, and keyword matching heuristics to assign ICD codes to clinical narratives~\citep{farkas2008automatic, scheurwegs2017data}. These systems often achieved high precision under narrowly defined contexts because medical experts designed rules explicitly targeting common or well-understood diagnoses. Nonetheless, maintaining rule-based systems for extensive code sets was difficult; an entire rule set would often need re-engineering whenever revisions to ICD or local documentation practices emerged~\citep{scheurwegs2017data}.

Parallel to rule-based approaches, researchers applied supervised classifiers such as logistic regression, Decision Trees, or Support Vector Machines (SVMs) trained on bag-of-words or TF-IDF representations of clinical notes~\citep{perotte2014diagnosis}. Although these methods sometimes benefited from straightforward interpretability (e.g., top weighted words for each code), they struggled with the extremely large and imbalanced label space, as ICD-9-CM and ICD-10-CM collectively encompass thousands of possible codes. Traditional feature engineering often could not capture the subtle semantic patterns indicative of specific diagnoses, particularly when these patterns were distributed across long or unstructured notes. Moreover, such approaches did not perform well on rare codes, which had scant training examples~\citep{duque2021knowledge}. Despite these limitations, these early systems established important baselines and indicated that integrating domain knowledge (e.g., simple dictionaries of medical terms) could enhance precision.

\subsubsection{Limitations and Motivation for Deep Learning}

Although traditional ML solutions revealed the potential of automated ICD coding, they were often constrained by high-dimensional data, sparse features, and a lack of capacity to fully capture clinical context. The low coverage of rare codes and the complex textual nature of clinical documentation prompted researchers to investigate more robust methods that did not rely heavily on labor-intensive rule sets or shallow textual features. The subsequent shift toward deep learning was motivated by these challenges, as neural architectures offered new pathways for automatically learning hierarchical textual representations with minimal handcrafting~\citep{lecun2015deep}.

\subsection{Deep Learning for Automated ICD Coding}\label{sec:p1:deep-learning}

\subsubsection{Convolutional and Recurrent Neural Network Approaches}

Deep learning opened up new avenues for capturing semantic and syntactic nuances in clinical text. Convolutional Neural Networks (CNNs) were among the earliest neural architectures to be tested, owing to their effectiveness in extracting local n-gram features~\citep{kim2014convolutional}. In the context of ICD coding, Mullenbach et al.~\citep{mullenbach2018explainable} introduced the Convolutional Attention for Multi-Label classification (CAML) model. CAML employed a convolutional encoder to represent textual segments, followed by a label-specific attention mechanism that highlighted the most relevant tokens for each code prediction. The increased interpretability of this attention mechanism was particularly appealing for healthcare applications, where transparency is critical:

\begin{equation}\label{eq:p1:caml-attention}
\alpha^{(l)} = \mathrm{softmax}(W^{(l)} h + b^{(l)}),
\end{equation}

where \(h\) denotes the hidden representations from the convolutional layers, and \(W^{(l)}\), \(b^{(l)}\) are label-specific parameters.

Subsequent CNN-based models explored deeper architectures with residual connections and multi-filter designs. For instance, Li and Yu~\citep{li2020multi} proposed the Multi-Filter Residual Convolutional Neural Network (MultiResCNN), which combines different kernel sizes and residual pathways to capture diverse syntactic and semantic patterns in clinical text. These refinements demonstrated measurable gains in F1-scores on benchmarks like MIMIC-III and other hospital datasets.

Recurrent Neural Networks (RNNs) offered an alternative strategy by encoding sequences of tokens in a step-by-step manner. Long Short-Term Memory (LSTM) networks~\citep{hochreiter1997long} and Gated Recurrent Units (GRUs) were particularly suited for capturing longer-range dependencies in clinical notes. Vu et al.~\citep{vu2020label, vu2020ijcai} employed a bidirectional LSTM with a label attention mechanism, where each label learned distinct attention weights over the text. This enabled more targeted coding decisions, accounting for the context needed by each diagnosis or procedure label. Although RNN models captured sequential information better than many CNN designs, they remained prone to vanishing gradients on very long documents and could be computationally expensive. 

By the late 2010s, CNN- and RNN-based models had markedly outperformed traditional ML baselines~\citep{perotte2014diagnosis}, demonstrating that neural networks could automatically learn powerful text features. Nevertheless, these architectures often required careful truncation or sampling of long notes, and they still struggled under conditions of data imbalance or when handling thousands of rarely seen ICD codes~\citep{wrenn2010quantifying}.

\subsection{Transformer-Based Architectures and Extensions}\label{sec:p1:transformers}

\subsubsection{Standard Transformers and Clinical Adaptations}

Transformer-based models, especially BERT~\citep{devlin2019bert}, revolutionized NLP by replacing the sequential encoding of RNNs with self-attention over entire sequences. This design facilitated more direct capture of long-range dependencies. However, standard BERT has a typical input limit of 512 tokens, which poses a challenge in clinical coding, where notes can span thousands of tokens~\citep{gao2021limitations}. Early attempts to integrate BERT into ICD coding often truncated notes, losing critical data. 

Huang et al.~\citep{huang2022plm} introduced the PLM-ICD model, which segments each discharge summary into manageable chunks and processes them with a pre-trained BERT. A label attention mechanism then aggregates chunk representations that are relevant to specific codes. Although such approaches allowed BERT-based solutions to outperform many CNN and RNN baselines on micro-F1 scores, truncated documents risked overlooking important evidence for certain diagnoses~\citep{dong2022automated}.

Heo et al.~\citep{heo2022medical} addressed this long-text issue by proposing the Document-to-Sequence BERT (D2SBERT) framework. It divides the full document into multiple fixed-length segments and processes each segment independently through BERT. A sequence-level attention mechanism then merges these representations, capturing essential information across the entire document. This method outperformed various CNN-based benchmarks on MIMIC-III, demonstrating that transformer-based models can effectively handle extensive notes through segmenting and hierarchical attention.

\subsubsection{Modern BERT-Based Encoders for ICD Coding}

Researchers also investigated architectural modifications and domain-specific pretraining to accommodate larger input lengths, improve inference speeds, and address the domain-specific vocabularies of clinical notes. Several recent works introduce novel BERT-like encoders:

\begin{itemize}
    \item \textbf{MosaicBERT}~\citep{portes2023mosaicbert} and \textbf{CrammingBERT}~\citep{geiping2023cramming}: Optimize training efficiency, allowing them to achieve comparable BERT performance with shorter training times. They do not fundamentally resolve the difficulties associated with extremely long narrative texts but reduce overall resource demands.
    \item \textbf{NomicBERT}~\citep{nussbaum2024nomicbert} and \textbf{GTE-en-MLM}~\citep{zhang2024gte}: Extend the maximum input length from 2k to 8k tokens and introduce advanced unpadding methods. These models exhibit modest performance improvements on retrieval tasks and partially address memory constraints by employing sparser attention patterns.
    \item \textbf{ModernBERT}~\citep{warner2024modernbert}: Incorporates architectural innovations such as alternating local-global attention, Rotary Positional Embeddings (RoPE), and label-aware attention modules. ModernBERT consistently outperforms prior baselines on both short and long ICD coding tasks, while reducing inference memory usage up to 8k tokens. Its hardware-aware design exemplifies how targeted architectural choices can mitigate computational overhead.
\end{itemize}

Despite progress, BERT-like transformers continue to encounter difficulties with extreme class imbalance, especially for codes encountered infrequently. Hardware constraints remain challenging when notes exceed thousands of tokens. Nonetheless, these tailored encoder designs underscore that carefully balancing domain knowledge, input length, and memory efficiency can produce state-of-the-art clinical coding models.

\subsection{Curriculum Learning and Hierarchical Knowledge Integration}\label{sec:p1:curriculum-learning}

The hierarchical structure of ICD codes naturally suggests a curriculum learning approach, where models are initially exposed to broad categories before proceeding to more granular ones. Bengio et al.~\citep{bengio2009curriculum} formally introduced curriculum learning as a strategy to improve model performance by scheduling examples from simple to difficult.

Ren et al.~\citep{ren2022hicu} applied a hierarchical curriculum (HiCu) by organizing ICD codes into tree-like levels. During training, they gradually refined predictions from higher-level ICD groupings (e.g., ``diseases of the circulatory system'') down to more specific subcodes. A hyperbolic embedding mechanism~\citep{nickel2017poincare} transferred knowledge from parent-level predictions to child-level nodes, enabling the model to leverage relationships among ICD codes. This approach mitigated some aspects of data imbalance, demonstrated by improved macro-F1 and macro-AUC scores, particularly for infrequent codes. Consequently, curriculum learning has emerged as a complementary strategy to neural architectures, aligning naturally with the structured taxonomies in ICD-9-CM, ICD-10-CM, and ICD-11.

\subsection{Challenges in Automated ICD Coding}\label{sec:p1:challenges}

\subsubsection{Data Imbalance and Rare Codes}

Automated ICD coding is inherently an extreme multi-label classification problem: a single patient note may contain multiple relevant codes among tens of thousands of possible labels. The distribution of ICD codes is typically long-tailed, where common chronic conditions such as hypertension or diabetes occur frequently, while countless rare conditions appear only a handful of times in training data~\citep{rios2018few}. In the MIMIC-III dataset, over 50\% of codes never appear or occur fewer than ten times~\citep{johnson2016mimic}, making it challenging for models to learn robust representations for infrequent diagnoses.

This imbalance leads to a disparity in model performance: systems that do well on common labels often fail to recall rare codes. Edin et al.~\citep{edin2023automated} conducted a replicability study that exposed how certain train-test splits and suboptimal macro F1 calculations led to inflated performances for rare-code prediction. Their work emphasized fairer evaluation and transparent reporting of model limitations. Similarly, Nguyen et al.~\citep{nguyen2023mimicivicd} showed that MIMIC-IV features a more extreme code distribution, since it merges ICD-9 and ICD-10 subsets. Handling this imbalance remains a top priority, as rare codes can indicate serious or complex conditions requiring accurate documentation and billing.

Curriculum learning strategies~\citep{ren2022hicu} partially address this issue by exploiting code hierarchies to share knowledge between parent and child labels. Additional techniques, such as oversampling and cost-sensitive learning, have also been tested, but they have been less effective at scale~\citep{scheurwegs2017data}. Consequently, novel approaches---particularly synthetic data augmentation---are garnering growing interest in mitigating the challenge of rare code prediction.

\subsubsection{Processing Long and Complex Clinical Documents}

Clinical narratives can be lengthy and may contain redundancies, inconsistencies, and a wide variety of clinical terminology. Early studies suggested that longer documents degrade model performance, but Edin et al.~\citep{edin2023automated} found the effect of length to be secondary to other factors such as code frequency. Nonetheless, managing long clinical notes remains computationally demanding.

Transformer-based models often face input length constraints (e.g., 512 or 1024 tokens), necessitating strategies such as truncation~\citep{gao2021limitations}, chunking~\citep{huang2022plm}, or hierarchical encoding~\citep{heo2022medical}. Although advanced architectures like Longformer or BigBird can theoretically handle up to 8k tokens or more, they are memory-intensive and can be slow to train. Document summarization and section-based modeling present additional strategies but risk losing crucial detail about specific diagnoses. Finding scalable, efficient ways to represent entire discharge summaries remains an important research focus for real-world deployments.

\subsubsection{Explainability and Interpretability}

Explainability is essential in clinical applications, where clinicians, auditors, and coding professionals must trust automated coding outputs~\citep{holzinger2017we}. Attention-based methods, such as CAML~\citep{mullenbach2018explainable} or LAAT~\citep{vu2020label}, highlight relevant text segments for each predicted code, providing intuitive evidence for model decisions. However, these attention weights only offer correlational explanations---further research is needed to establish more causal or structured forms of explainability.

Edin et al.~\citep{edin2024unsupervised} proposed an unsupervised approach that can achieve supervised-level explanation quality, suggesting that external knowledge or adversarial training may produce rationale annotations without costly manual labeling. Additionally, researchers have explored knowledge graphs and symbolic reasoning to justify individual code predictions~\citep{teng2020explainable}. Although these methods can yield greater transparency, fully interpretable models that seamlessly integrate domain knowledge remain an open challenge.

\subsubsection{Embedding Models and Clinical Semantic Search}

Embedding models underlie many of the text representation and retrieval tasks in healthcare. They transform sparse textual data into dense vectors that encode semantic relationships among words, sentences, or entire documents. Excoffier et al.~\citep{excoffier2024generalist} compared generalist transformers (trained on broad corpora) to specialized clinical embeddings for ICD semantic search, observing that generalist embeddings were often more robust to minor text rephrasings. This finding implies that domain-specific pretraining might not always be superior if it narrows the model's domain to patterns less frequent or less stylistically varied than the general language.

Nevertheless, domain-adapted models like ClinicalBERT, BioBERT, and PubMedBERT have achieved better performance in several medical NLP tasks, including ICD coding~\citep{gao2021limitations, ji2021bert}. The discrepancy in outcomes suggests that training set size, the nature of clinical text, and the distribution of code descriptions can all influence whether generalist or specialized embeddings hold the advantage. Moreover, specialized strategies such as code-synonym expansion~\citep{yuan2022codesynonyms} and concept-level attention highlight how embedding-based approaches can leverage existing biomedical ontologies (e.g., UMLS~\citep{bodenreider2004umls}, SNOMED CT~\citep{donnelly2006snomedct}, or cTAKES~\citep{savova2010ctakes}) to improve both performance and explainability.

\subsubsection{Knowledge-Based Methods and Symbolic Integration}

Knowledge-based methods aim to incorporate medical ontologies, taxonomies, and domain rules into deep learning architectures. These techniques can address shortcomings in purely data-driven models by providing explicit relationships between diseases or leveraging synonyms and hierarchical groupings~\citep{xie2019ehr, teng2020explainable}. Such integration is especially valuable in ICD coding, where label space is inherently hierarchical, and textual evidence for certain diagnoses may be sparse.

Ren et al.~\citep{ren2022hicu} demonstrated that hyperbolic embeddings capture hierarchical structure better than Euclidean embeddings, particularly for ICD codes. By representing codes in a hyperbolic space, parent-child relationships are learned with low distortion, and training proceeds from broader categories to finer levels. Similarly, Cao et al.~\citep{cao2020hypercore} proposed a hyperbolic co-graph representation that leveraged hierarchical information to enhance predictive performance. Knowledge-based approaches thus help unify symbolic reasoning with data-hungry neural models, thereby improving interpretability and filling gaps where training examples are limited.

\subsection{Synthetic Data Generation for Rare ICD-10-CM Codes}\label{sec:p1:synthetic-data}

\subsubsection{Motivation and Techniques}

The extreme class imbalance of ICD codes---where many rare diagnoses have few or no training examples---has motivated synthetic data generation as a complementary strategy to curriculum learning and data augmentation. Large Language Models (LLMs), such as GPT-3.5, GPT-4, or fine-tuned LLaMA, are increasingly employed to create realistic clinical narratives for underrepresented diagnoses~\citep{kumichev2024medsyn, falis2024gpt35}. These synthetic notes preserve confidentiality while expanding the variety of text samples available for training models.

Falis et al.~\citep{falis2024gpt35} examined the feasibility of prompting GPT-3.5 with specific ICD-10-CM codes and descriptions to generate discharge summaries. Adding these synthetic notes to real training data produced modest performance improvements on rare codes and helped the model avoid out-of-family misclassifications. However, these gains came at the cost of potentially overfitting to the style of generated text, which might be less diverse than real clinical documentation.

Other augmentation methods include synonym replacement, back-translation, or adversarial generation of latent feature vectors, rather than raw text. Although these approaches are less computationally demanding than prompting large LLMs, they often provide smaller increases in coverage for rare labels. Hybrid strategies that combine knowledge-graph constraints with generative models have also been explored, ensuring that synthetic samples remain clinically plausible~\citep{falis2022horsestozebras}.

\subsubsection{Ethical and Practical Considerations}

Synthetic data introduce potential advantages for privacy protection, as real patient identifiers can be omitted or obfuscated. Nevertheless, generative models might inadvertently reproduce sensitive or identifiable fragments if they encountered them during training. Ensuring high-fidelity clinical plausibility is also essential; simplistic or repetitive patterns within synthetic samples could bias the model toward unrepresentative language~\citep{pollard2018eicu}.

In practice, synthetic notes appear to be most beneficial in few-shot or one-shot contexts, where the model has a handful of real examples for a rare code and gains additional variety through generation~\citep{rios2018few}. Full replacement of real data is typically detrimental. As a result, synthetic data remain a supplement rather than a primary information source, especially for regulating or auditing high-stakes billing environments.

\subsection{Benchmarks and Evaluation Metrics}\label{sec:p1:benchmarks}

\subsubsection{Public Datasets and Benchmarks}

The research community has benefited significantly from openly available clinical datasets such as MIMIC-III~\citep{johnson2016mimic} and MIMIC-IV~\citep{johnson2023mimicivscidata}. MIMIC-III focuses on ICD-9-CM codes, while MIMIC-IV provides a larger number of hospital admissions and reflects the transition to ICD-10-CM coding~\citep{nguyen2023mimicivicd}. Both datasets contain lengthy discharge summaries with manually assigned ICD codes, making them ideal for evaluating multi-label classification performance. The eICU Collaborative Research Database~\citep{pollard2018eicu} supplies additional notes and structured data across multiple hospital systems, enabling analyses of model generalizability.

Benchmark initiatives standardize tasks, train-dev-test splits, and evaluation pipelines. This fosters comparability among different architectures (e.g., CNNs vs. BERT vs. hierarchical models) and clarifies progress in rare-code prediction. Replication studies have demonstrated that small variations in preprocessing or code filtering can materially affect reported metrics, emphasizing the need for consistent benchmark protocols~\citep{edin2023automated}.

\subsubsection{Common Evaluation Metrics}

Because automated ICD coding is a multi-label task, precision and recall must be balanced to capture different aspects of performance:

\begin{itemize}
    \item \textbf{Micro-F1}: Aggregates all predicted and true labels across the dataset, giving higher weight to frequently occurring codes. This is typically higher than macro-F1 and often used to illustrate a model's overall performance on the bulk of common codes.
    \item \textbf{Macro-F1}: Computes the F1-score per code and then averages across all codes, giving equal importance to rare and common labels. Macro-F1 is sensitive to performance on the long tail of infrequent codes and is usually much lower.
    \item \textbf{Micro/Macro AUC}: Measures the area under the ROC curve, either pooling predictions across codes (micro) or taking a simple average of per-code AUC (macro). High AUC often indicates that, with careful threshold tuning, the model could achieve good recall and precision.
    \item \textbf{Precision@K}: Assesses how many of the top \(K\) predicted codes match the ground truth. This metric is sometimes used in practical coding assistance tools to limit the suggestions displayed to human coders.
    \item \textbf{Hierarchical Metrics}: Account for partial correctness if a model predicts a closely related code from the same family, as opposed to a completely unrelated code. Such metrics better reflect real-world coding tasks, where distinguishing the correct subtype can be secondary to identifying the right disease category.
\end{itemize}

Error analyses often examine whether incorrect predictions are at least within the correct organ system or parent node in the ICD hierarchy. Some studies also evaluate interpretability by checking if highlight-based explanations align with key phrases annotated by clinical experts~\citep{mullenbach2018explainable}. These deeper analyses provide a more nuanced picture of model behavior than a single F1 score.

\subsection{Summary of ICD Coding Literature}\label{sec:p1:icd-summary}

Automated ICD coding presents a set of distinct technical and organizational challenges. Initial studies employed rule-based logic and feature-engineered machine learning but were constrained by scalability and the evolving nature of clinical vocabularies~\citep{scheurwegs2017data}. Deep neural architectures, particularly CNNs, RNNs, and transformers, provided significant accuracy improvements by learning hierarchical text representations without extensive manual engineering~\citep{mullenbach2018explainable, li2020multi, heo2022medical}.

Transformer-based models and newer BERT-like encoders excel at capturing complex linguistic features in discharge summaries. However, they require careful handling of long documents, face persistent issues with rare diagnoses, and demand large computational resources. Curriculum learning, hierarchical embeddings, and synthetic data generation offer partial remedies, enhancing performance on uncommon codes and improving interpretability through structured knowledge infusion~\citep{ren2022hicu, falis2024gpt35}.

Recent benchmarking efforts clarify that future progress lies in robust domain-specific pretraining, advanced handling of extreme multi-label imbalances, and deeper integration of medical ontologies. Additionally, evolving demands---such as ICD-11 coding---underline the importance of flexible systems that adapt to changing taxonomies. The development of interactive, interpretable coding assistants, combined with rigorous evaluation and continuous learning, is expected to guide the field toward reliable real-world implementation in clinical and administrative workflows.

%% ============================================================
%% PART II — SA-AKI Mortality Prediction (Current Report, preserved verbatim)
%% ============================================================
\section{Part II: SA-AKI Mortality Prediction}\label{sec:saaki-lit}

\subsection{Clinical definitions and epidemiological context}
Sepsis-associated acute kidney injury (SA--AKI) has been characterized as the concurrence of sepsis and acute kidney injury with persistently poor outcomes. Sepsis has been defined by the Sepsis--3 consensus as life-threatening organ dysfunction arising from a dysregulated host response to infection, and the construct has been operationalized as suspected or documented infection together with an acute rise in organ failure, commonly a SOFA increase of at least two points \citep{Singer_2016}. Acute kidney injury has been defined by KDIGO through serum creatinine and urine output thresholds, in particular an absolute creatinine increase of at least $0.3\,\mathrm{mg/dL}$ within forty-eight hours, or at least $1.5\times$ baseline within seven days, or urine output of $0.5\,\mathrm{mL/kg/h}$ or less for six hours \citep{KDIGO_2012}. Many studies have labeled SA--AKI when AKI develops within a short interval after sepsis onset, frequently within seven days, to attribute renal dysfunction to the septic process rather than to competing causes \citep{Peerapornratana_2019,Poston_Koyner_2019}. The shift toward Sepsis--3 and KDIGO has reduced heterogeneity relative to older criteria and has improved comparability across cohorts \citep{Singer_2016,KDIGO_2012}.

The burden of SA--AKI in intensive care units is substantial. AKI develops in roughly fifty to sixty percent of septic ICU patients, and sepsis accounts for about forty to seventy percent of AKI in that setting \citep{Peerapornratana_2019,Poston_Koyner_2019}. When AKI complicates sepsis, marked increases in in-hospital mortality follow, with earlier analyses suggesting a six to eight fold rise compared with sepsis without AKI and mortality rates commonly between forty and sixty percent \citep{Poston_Koyner_2019,Peerapornratana_2019}. Prolonged ICU stays and an elevated risk of subsequent chronic kidney disease have also been recorded among survivors \citep{Poston_Koyner_2019,Peerapornratana_2019}. These epidemiological patterns have motivated the development of early risk stratification tools.

\subsection{Conventional severity indices and early warning systems}
General ICU indices such as SOFA, SAPS~II, and APACHE~II have been evaluated for SA--AKI outcomes and tend to show only modest discrimination for short-term mortality or kidney-specific endpoints in this subgroup \citep{Luo_2022,Li_2023_SciRep,Qu_2024}. Classical regression frameworks, including logistic regression and Cox models, have been employed repeatedly, yet difficulties persist in capturing the nonlinear and interacting relationships among infectious physiology, hemodynamics, renal injury, and treatment exposures. Performance has often failed to transport reliably across hospitals as a result \citep{Kang_Yoon_2023,Luo_2022,Li_2023_SciRep}. Earlier variability in sepsis and AKI definitions compounded the challenge by complicating cohort selection and cross-study comparison prior to widespread adoption of Sepsis--3 and KDIGO \citep{Singer_2016,KDIGO_2012}.

\subsection{Datasets and cohort construction in SA--AKI modeling}
Large electronic health record repositories have supported the development of prognostic models for SA--AKI, with MIMIC--IV serving as the most common source for cohort construction and feature derivation \citep{Johnson_2023_MIMIC}. Studies typically restrict analysis to the first ICU admission per patient, exclude patients younger than eighteen years, and remove ICU stays shorter than forty-eight hours. A central design decision involves the temporal window for feature construction. Fixed early windows, for example the first twenty-four or forty-eight hours after admission or sepsis recognition, support timeliness, whereas whole-stay aggregation from admission to discharge has also appeared frequently and creates a risk that future information relative to an intended decision point enters the model inputs.

\subsection{Feature engineering, selection, and leakage risks}
The handling of high-dimensional clinical variables has relied on screening by correlation or variance inflation factor, wrapper methods such as recursive feature elimination, and embedded strategies that exploit model-based importance estimates. It is common to find feature selection steps that were informed by the full dataset. When selection proceeds in this manner, the final predictor set is influenced by outcomes from records that are later evaluated, and optimism in reported discrimination is encouraged \citep{Kang_Yoon_2023}. Claims that VIF-based pruning enhances stability and generalizability have also appeared without pre--post diagnostics or external tests, and the relevance of VIF is limited for classes of models that tolerate correlated predictors. In several studies, length of stay was included among features for mortality prediction. Because time to event is strongly associated with whether death occurred during the admission, that inclusion acts as a conduit for post-baseline information and dominates other covariates during training. Redundancy among variables has also been noted. For example, total urine output and average urine output become scalar multiples under a fixed observation horizon, and diverging importance profiles for those variables suggest that non-fixed or whole-stay windows were used, which reduces alignment with early prediction. Dependencies between creatinine and derived eGFR similarly complicate interpretation when both are included without a clear rationale.

\subsection{Missing data mechanisms and imputation choices}
The mechanisms that govern missingness have rarely been examined in depth in this literature. Assumptions that data were missing completely at random or missing at random are often stated without empirical assessment, and simple percentage-based rules trigger imputation. Mean imputation for continuous variables remains a common choice even for biomarkers that show marked skew, such as serum creatinine, which is right skewed and well approximated by a log-normal distribution \citep{Haeckel_2010}. Multiple imputation by chained equations and random forest variants is widely used, and in many cases those imputation models are fitted across the entire dataset prior to train--test separation. Because chained equations use all observed variables to impute each target variable, global fitting allows information from records later treated as held out to influence imputed values, and optimistic estimates of model performance follow.

\subsection{Class imbalance, evaluation strategy, and thresholding}
Mortality among SA--AKI cohorts is imbalanced, and class imbalance is most often handled by oversampling techniques, including SMOTE and ADASYN \citep{Chawla_2002,He_2008}. The two methods have occasionally been applied together even though ADASYN is itself a variant that emphasizes synthesis in hard regions of the feature space. SMOTE may generate synthetic points in overlapped or noisy areas, and performance can degrade in high-dimensional settings unless additional filtering is applied \citep{Blagus_2013,Saez_2015}. Concentration of synthesis near low-density or borderline regions by ADASYN can shift decision boundaries around outliers and may harm stability under severe imbalance or distribution shift \citep{Krawczyk_2016}. Comparisons with alternative imbalance strategies, for instance through intrinsic class weighting or hybrid sampling, have not always been provided. Evaluation strategies have often relied on a single split of the data into training and test sets without a separate validation split, and explicit calibration analyses and principled threshold selection have been omitted in several reports.

\subsection{Model classes, tuning, calibration, and interpretability}
Tree-based ensembles, notably extreme gradient boosting, have been adopted widely in this domain to represent nonlinearities and interactions among clinical variables, and stronger internal discrimination than baseline ICU scores is often reported \citep{Chen_Guestrin_2016,Luo_2022,Li_2023_SciRep,Qu_2024,Wang_2024,Chen_2025,Le_2025}. Hyperparameter tuning is usually performed through grid search cross validation, and claims of optimality appear despite the presence of complex, nonconvex search spaces in which only local minima are guaranteed by such procedures. Calibration is reported inconsistently, and explanations based on Shapley values are used to describe feature contributions, although stability of such explanations is not always explored. In several comparisons, confidence intervals for ensembles and penalized linear models overlap, and differences attributed to model class are entangled with choices about temporal windows, selection scope, and imputation design.

\subsection{External validity, transportability, and deployment timing}
Generalization beyond the MIMIC--IV development site remains a consistent concern. MIMIC--IV reflects a single tertiary-care center in Boston with local care patterns and population structure that are not broadly representative \citep{Johnson_2023_MIMIC}. External evaluations of ICU prediction models show loss of accuracy when transported across hospitals, pointing to center-specific differences that matter for performance \citep{Rockenschaub_2024}. Assertions about applicability to diverse populations therefore require cautious interpretation when external testing is absent. In addition, the timing of features relative to the intended decision point is often unspecified, and whole-stay aggregation creates ambiguity about whether a model could power actionable alerts at twenty-four or forty-eight hours.

\subsection{Detailed appraisal of representative studies}

\subsubsection{Chen et al.\ (2025): SA--AKI mortality in MIMIC--IV}
A recent contribution by Chen et al.\ reported an AUROC of 0.878 for mortality prediction in a SA--AKI cohort drawn from MIMIC--IV, and the result was presented as an improvement over the AUROC 0.794 attributed to Li et al.\ \citep{Chen_2025,Li_2023_SciRep}. The cohort retained only the first ICU admission per patient and excluded patients younger than eighteen years and ICU stays shorter than forty-eight hours. The modeling pipeline combined VIF-based filtering with recursive feature elimination to produce a set of twenty-four predictors, after which several variables were reintroduced as expert inputs without quantitative justification. Because selection was described as having been performed on the entire dataset, the final predictor set was influenced by outcomes from records later treated as held out and the risk of information leakage was created \citep{Kang_Yoon_2023}. A claim of generalizability across diverse populations was made, yet MIMIC--IV represents a single center and patient mix, and performance drift across hospitals has been documented \citep{Johnson_2023_MIMIC,Rockenschaub_2024}. The introduction criticized Li et al.\ for relying on forty-four variables on the grounds that the set might be too narrow to capture the complexity of SA--AKI, although the final model by Chen et al.\ used an even smaller set of twenty-four features, which created an internal inconsistency in the narrative \citep{Chen_2025,Li_2023_SciRep}. Hyperparameters were tuned with grid search cross validation and described as optimal, although such procedures identify settings within a discrete grid and do not ensure global optimality in nonconvex spaces. Claims of timely, accurate, and actionable interventions were advanced, yet the temporal window for aggregating time-varying biomarkers was not specified. Because Sepsis--3 and KDIGO anchor observation to early horizons, an early extraction window would have been expected for actionable deployment; multiple signals suggested that whole-stay aggregation may have been used, which would introduce future information relative to early prediction. Missing values were imputed by means that were tied to percentage thresholds, and mean imputation was applied for variables with up to twenty percent missingness. This approach is not well aligned with skewed distributions such as serum creatinine, which has been described as right skewed and is well modeled by a log-normal distribution \citep{Haeckel_2010}. Missingness indicators were not created and the mechanism of missing data was assumed rather than studied. VIF filtering was described as enhancing stability and generalizability without pre–post diagnostics or external assessment and has limited relevance for ensembles that tolerate correlated predictors. The base estimator used within recursive feature elimination was not stated, and a linear estimator would impose a separability assumption during selection. Length of stay was included among predictors, and this inclusion introduced a strong proxy for time to event that is not available prospectively for early prediction. The simultaneous inclusion of total and average urine output suggested redundancy if a fixed window had been used, and differences in importance between the two variables implied that a non-fixed or whole-stay window was applied. The evaluation strategy used a seventy-five to twenty-five train–test split without a separate validation split, and explicit calibration and threshold selection were not described. Class imbalance was handled by applying SMOTE and ADASYN together, a combination that lacks a clear rationale and that may synthesize examples in overlapped or noisy regions or near outliers and thereby distort decision boundaries \citep{Chawla_2002,He_2008,Blagus_2013,Saez_2015,Krawczyk_2016}. Univariate comparisons between survivors and a group termed non-survivors were presented with t tests even though laboratory variables such as creatinine generally deviate from normality \citep{Haeckel_2010}. The reported confidence intervals for ensemble and linear models overlapped, with AUROC 0.878 (95\% CI 0.859 to 0.897) for extreme gradient boosting and AUROC 0.849 (95\% CI 0.824 to 0.873) for logistic regression, which supported the interpretation that any advantage might be smaller than suggested once temporal alignment and leakage are addressed \citep{Chen_2025}.

\subsubsection{Li et al.\ (2023): SA--AKI mortality in MIMIC--IV}
Li et al.\ assembled a SA--AKI cohort from MIMIC--IV by retaining only the first ICU admission and excluding ICU stays shorter than forty-eight hours. Features were extracted from the first twenty-four hours and were aggregated using maximum values, which provided a clear early window for model inputs \citep{Li_2023_SciRep}. Lasso-regularized logistic regression was used for feature selection, and selection was described across the full dataset rather than within development folds. Because selection then involved outcomes from records later used for evaluation, optimism in reported performance was encouraged. The linear framework also prioritized main effects and could miss interactions and subtle nonlinear patterns that tree ensembles tend to capture. Multiple imputation by chained equations with random forest implementations was applied globally across the dataset rather than learned on development data and then applied to evaluation data, and missingness indicators were not constructed. The data were split into training and validation in an eighty to twenty ratio, and a discrimination gap favoring extreme gradient boosting over logistic regression was reported with AUROC 0.794 versus 0.730, although the magnitude of this gap remained difficult to interpret in light of the selection and imputation scope \citep{Li_2023_SciRep}.

\subsubsection{Hu et al.\ (2021): survival analysis in MIMIC--III}
A retrospective analysis by Hu et al.\ used MIMIC--III with smaller sample size and older coding structure to study SA--AKI outcomes \citep{Hu_2021}. The cohort retained only the first ICU admission and excluded patients younger than eighteen years and ICU stays shorter than forty-eight hours, and rows with more than five percent missing variables were removed. Lasso regression and Cox regression were used, and a C index of 0.72 was reported on a validation set. Biomarkers were aggregated by medians across the entire admission, therefore features included information that would not be available at an early decision point. Features with more than twenty percent missingness were dropped without investigation of mechanisms, and multiple imputation by chained equations was performed on the full dataset prior to splitting, which created leakage. The test set contained one hundred and two rows, and reliability of performance estimates was limited by that size. Ethnicity was used as a predictor, which raised concerns about representativeness and fairness for transportability. Survival models at seven, fourteen, and twenty-eight days yielded AUCs of 0.77, 0.73, and 0.71 respectively.

\subsubsection{Roknaldin et al.\ (2024): sepsis to AKI prediction in MIMIC--III}
A recent preprint constructed models that predict AKI among septic ICU patients using MIMIC--III \citep{Roknaldin_2024_arXiv}. Adults aged eighteen to eighty-nine years with sepsis by ICD--9 codes were included after excluding multiple admissions, ICU stays shorter than forty-eight hours, and records with more than twenty percent missingness. Features from the first twenty-four hours covered demographics, comorbidities, vitals, laboratory results, and early interventions. Multiple imputation was used for missing values, a correlation filter retained twenty-three predictors, and class imbalance was addressed with SMOTE. Six algorithms were compared, and logistic regression achieved the strongest internal results with test AUC 0.887 and validation AUC 0.857 together with a Brier score near 0.13 and acceptable isotonic calibration. Top predictors included urine output, bilirubin measures, blood urea nitrogen, minimum eGFR, weight, and indicators of vasopressor use and mechanical ventilation. The transparency of the feature list and the early window supported interpretability, while several design choices encouraged optimism. Feature selection and SMOTE were described prior to the final split into development and evaluation, so that the set of predictors and synthetic examples were influenced by records that later appeared in evaluation. Multiple imputation also appeared to have been fitted globally. A fixed correlation threshold risked removal of variables that are weakly correlated but jointly informative, and the inclusion of minimum and maximum eGFR alongside creatinine introduced redundancy because eGFR is derived from creatinine and demographics. For descriptive group comparisons, t tests were used for markers that are typically skewed, and the allocation of half the data to validation and one tenth to testing left a small test set. Calibration procedures were not tied to a recalibration step that excluded the test cohort, which limited the strength of calibration claims on that set.

\subsubsection{Chaudhary et al.\ (2020): phenotype discovery in MIMIC--III}
An unsupervised study in \emph{CJASN} sought to identify SA--AKI phenotypes rather than predict mortality directly \citep{Chaudhary_2020_CJASN}. An autoencoder was used to learn representations that were then clustered into three subgroups with distinct twenty-eight day mortality and dialysis rates. Differences were driven by bilirubin, lactate, transaminases, and inflammatory counts, and the highest risk subgroup showed markedly higher dialysis use and death. The analysis relied on whole-stay trajectories of laboratory and vital signs, which embeds information not available at early decision times. The sepsis cohort was constructed by codes, which can misalign the timing of the septic insult, and external validation of cluster stability was not performed. The exploratory links between clusters and treatments were not framed as causal and would require dedicated designs to support treatment guidance.

\subsubsection{Lin et al.\ (2025): progression to chronic kidney disease in MIMIC--IV}
A study in \emph{Informatics in Medicine Unlocked} examined progression from SA--AKI to incident chronic kidney disease using MIMIC--IV \citep{Lin_2025_IMU}. Extreme gradient boosting with SHAP explanations was trained on features collected during the first twenty-four hours, and high internal discrimination was reported. Thresholds for age, maximum creatinine, minimum hemoglobin, and lactate were described, and a feature stability analysis across AKI stages was presented. Although the early window and the inclusion of decision-curve analysis aided interpretability, variable screening by Lasso and stepwise logistic regression prior to final modeling, if performed across the full dataset, would have allowed outcomes from evaluation records to influence the screened set. Outcome assessment relied on follow up up to five years, and selection by loss to follow up as well as center-specific practice patterns could have influenced estimates. Thresholds based on SHAP dependence plots have known sensitivity to sampling variation, and replication would be expected for reliable adoption.

\subsubsection{He et al.\ (2021): acute kidney disease after SA--AKI}
An article in \emph{Frontiers in Medicine} predicted subsequent acute kidney disease rather than in-hospital mortality \citep{He_2021_FrontMed}. A single-center development cohort was used, and external validation was undertaken in MIMIC--III. Recurrent neural networks with long short-term memory units were compared with a decision tree and logistic regression. The recurrent model achieved an internal AUROC of 1.00 and strong performance in the external cohort, and the change in nonrenal SOFA between day one and day three emerged as an important signal. The development cohort contained two hundred and nine patients, which is a modest sample for a deep architecture with more than two dozen features, and perfect discrimination in such data is difficult to reconcile with typical ICU noise and suggests that overfitting, simplified labels, or other forms of information leakage may have been present. The dynamic window implied by the day one to day three delta improves clinical timing, while the time of alert generation relative to clinical workflow was not specified. It was not always clear that imputation and preprocessing steps were confined to development data.

\subsubsection{Kantola et al.\ (2025): admission procalcitonin and risk of AKI}
A multicenter retrospective cohort across three hospitals evaluated admission procalcitonin for prediction of subsequent AKI and reported an odds ratio of approximately 1.01 per ng/mL with poor overall discrimination near ROC AUC 0.59 \citep{Kantola_2025}. The heterogeneity of practical cutoffs across prior work and the inconsistent relationships between procalcitonin levels and AKI severity were emphasized, and the conclusion favored limited standalone clinical utility for early risk stratification. Procalcitonin is not available in MIMIC--IV, which restricts replication in that database.

\subsection{Synthesis and implications for this thesis}
The literature reviewed above sketches a consistent picture in which window definition, the scope of data-dependent preprocessing, handling of missingness and imbalance, and the availability of external evidence play decisive roles in the credibility of reported scores. Studies that relied on whole-stay aggregation, on inclusion of length of stay or other post-baseline correlates, and on global feature selection or imputation tend to report stronger discrimination, yet those gains are paired with ambiguity about practical timing and with higher risk of optimism. When early windows are enforced clearly, performance remains competitive and the path to bedside use becomes clearer. Several reports present overlapping confidence intervals between ensembles and penalized linear models, which suggests that differences in pipeline design often overshadow differences in the learning algorithms themselves \citep{Kang_Yoon_2023}. Generalization beyond a single tertiary-care center remains a key uncertainty, and drift across hospitals is documented \citep{Johnson_2023_MIMIC,Rockenschaub_2024}. In this thesis, these observations are treated as design constraints for methodological choices, and the subsequent chapters are organized to respond to the gaps identified here while keeping the emphasis on temporally anchored modeling, leakage-aware preprocessing, calibrated evaluation, and transparency about what information is and is not available at the moment when a clinical decision is intended to be supported.

%% ============================================================
%% SYNTHESIS — Tying both bodies of literature together (NEW)
%% ============================================================
\section{Synthesis: Machine Learning for Clinical Decision Support Using MIMIC-IV}\label{sec:lit-synthesis}

The two bodies of literature reviewed in this chapter---automated ICD coding and SA-AKI mortality prediction---address distinct clinical tasks yet share a common data foundation in the MIMIC-IV database and converge on several methodological themes that define the current frontier of machine learning for clinical decision support.

First, both domains contend with severe class imbalance. In ICD coding, the long-tailed distribution of diagnostic codes means that rare but clinically significant conditions receive inadequate representation during training~\citep{rios2018few, nguyen2023mimicivicd}. In SA-AKI mortality prediction, the minority class of non-survivors is substantially outnumbered by survivors, and naive oversampling strategies risk distorting decision boundaries~\citep{Chawla_2002, Krawczyk_2016}. The parallel challenge underscores the need for principled imbalance handling---whether through curriculum learning and synthetic data generation for ICD coding, or through intrinsic class weighting and calibrated thresholding for mortality prediction.

Second, both literatures highlight the tension between model complexity and interpretability. ICD coding research has pursued attention-based explanations~\citep{mullenbach2018explainable} and knowledge-graph integration~\citep{teng2020explainable} to make predictions transparent to clinical coders. SA-AKI prediction studies have relied on SHAP-based feature attribution and decision-curve analysis to demonstrate clinical face validity. In both cases, the goal is to produce models whose outputs clinicians can trust and act upon, a prerequisite for any clinical decision support system.

Third, the MIMIC-IV database serves as the shared empirical substrate. For ICD coding, MIMIC-IV provides discharge summaries with manually assigned ICD-9 and ICD-10 codes, enabling evaluation of multi-label classification at scale~\citep{johnson2023mimicivscidata}. For SA-AKI prediction, the same database supplies structured clinical variables, time-stamped laboratory results, and outcome labels that support cohort construction and temporal feature engineering~\citep{Johnson_2023_MIMIC}. This shared foundation creates a natural opportunity to study both problems within a unified analytical framework, leveraging common data extraction pipelines and ensuring consistency in patient populations.

Fourth, both domains face challenges of external validity. ICD coding models trained on MIMIC-III or MIMIC-IV may not generalize to institutions with different documentation practices or coding conventions~\citep{edin2023automated}. SA-AKI prediction models developed on MIMIC-IV reflect a single tertiary-care center and show documented performance drift when transported to other hospitals~\citep{Rockenschaub_2024}. Addressing this limitation requires either multi-site validation or careful acknowledgment of the boundaries within which model predictions remain reliable.

Finally, both research threads emphasize the importance of methodological rigor in preprocessing. In ICD coding, inconsistent benchmark protocols and preprocessing choices materially affect reported metrics~\citep{edin2023automated}. In SA-AKI prediction, global feature selection and imputation prior to train-test separation introduce information leakage that inflates discrimination estimates~\citep{Kang_Yoon_2023}. The lessons from both literatures converge on a shared imperative: transparent, reproducible pipelines with clearly defined temporal boundaries and leakage-aware design.

This thesis addresses both problems under a unified clinical decision support framing. Problem~1 (automated ICD coding for rare diseases) draws on the NLP and deep learning literature reviewed in Part~I, while Problem~2 (SA-AKI early mortality prediction) builds on the clinical modeling literature reviewed in Part~II. The subsequent chapters present the problem formulations, methodologies, datasets, and results for each problem, with MIMIC-IV as the common data foundation throughout.
